\documentclass[12pt,letterpaper]{article}
\usepackage{amsmath,amssymb,amsthm}
\usepackage{newtxtext,newtxmath}
\usepackage[margin=1in]{geometry}
\usepackage{microtype}
\usepackage{enumitem}
\usepackage{longtable,booktabs,array}
\usepackage[colorlinks=true,urlcolor=blue,linkcolor=blue]{hyperref}
\setlength{\parindent}{15pt}
\setlength{\parskip}{3pt}
\setlist[enumerate]{leftmargin=*,itemsep=4pt}
\newcommand{\F}{F}
\newcommand{\OF}{\mathcal O_F}
\newcommand{\E}{E}
\newcommand{\G}{\mathrm{PGL}_3}
\newcommand{\Hone}{\mathrm{SL}_2}
\newcommand{\Htwo}{\mathrm{PGL}_2}
\newcommand{\Fix}{\operatorname{Fix}}
\newcommand{\Sym}{\operatorname{Sym}}
\newcommand{\tr}{\operatorname{tr}}
\newcommand{\ord}{\operatorname{ord}}
\newcommand{\Nm}{\operatorname{N}}
\newcommand{\cl}{\operatorname{cl}}
\newcommand{\pr}{\mathbb P}
\title{Referee Report on\\[4pt]
\emph{On twisted orbital integral identities for $PGL(3)$ over a $p$-adic field}}
\author{Independent mathematical and editorial assessment}
\date{23 November 2026}
\begin{document}
\maketitle

\begin{abstract}
This report assesses the correctness and readability of the six-part
\LaTeX{} manuscript \texttt{adjoint\_lift0.tex}--\texttt{adjoint\_lift5.tex},
with \texttt{adjoint\_lift\_pics.tex} considered as the associated figure
source. The main published result is a geometric comparison of orbital
integrals for the unit Hecke function, under the odd residual-characteristic
hypothesis. The argument has a coherent overall architecture, but the supplied
text contains several demonstrable inconsistencies and a small number of
claims that are false literally. Some predate the 2003 re-typesetting; at
least two mathematical transcription mistakes can be identified by direct
comparison with the 1990 journal publication. The recommendation is a
\emph{corrected historical edition with a substantive errata appendix}, not
an unqualified republication in the present form. This is a targeted
mathematical audit, not a complete independent reconstruction of every
building-counting proof in Langlands' unpublished notes.
\end{abstract}

\section{Recommendation and scope}

\noindent\textbf{Recommendation: major corrections and explicit verification
before treating this version as a mathematically authoritative edition.}
This recommendation concerns the present text, and is not a claim that the
published fundamental-lemma result is false. The strategy is persuasive: the
split case is handled by Satake transforms, the nonsplit case is reduced to
fixed-point counts on Bruhat--Tits buildings, and the distinct geometric
contributions are compared in Sections 4--5. The figures are helpful for
visualizing characteristic leaves and their segments. However, several
formulas defining the objects to be compared are not correct as written.

The audit has three levels: (i) direct algebraic or logical checks;
(ii) checks of the dependency and normalization chain in Sections 3--5;
and (iii) spot comparison with the published article, Joyner
\cite{Joyner1990}. I have not independently checked the unpublished
Langlands manuscript or all of the implicit building combinatorics; an
unmentioned step is \emph{not} thereby certified correct. Throughout,
references such as ``(3.18)'' are to the paper's own equation numbering.
Line numbers refer to the individual uploaded source files, not to the
concatenated edition.

\section{High-priority mathematical findings}

\subsection{The symmetric-power representation is mislabeled and the basis
for (1.3) is wrong}
\label{issue:adjoint}

\textbf{Location:} \texttt{adjoint\_lift1.tex}, lines 49--93;
\textbf{priority: high}. The map
\[
r:\mathrm{GL}_2(\mathbb C)\longrightarrow\mathrm{GL}_{2n+1}(\mathbb C)
\]
is described as
$\Sym^{\otimes 2n}\otimes(\det)^{-n}$. The tensor power of a
2-dimensional standard representation has dimension $2^{2n}$, not
$2n+1$. The representation with the specified target dimension is
\[
\boxed{\Sym^{2n}(\mathrm{Std})\otimes(\det)^{-n}}.
\]
The phrase ``symmetric $n$th power lifting'' likewise needs qualification:
in this notation it is the symmetric \emph{$2n$th} power. This issue is
also visible in the 1990 typeset article, p.~1099 \cite{Joyner1990}.

There is a further, independently testable problem with the basis announced
immediately after matrix (1.3). The displayed middle matrix
$\bigl(\begin{smallmatrix}1&0\\-1&0\end{smallmatrix}\bigr)$
has trace $1$ and therefore does not belong to
$\mathfrak{sl}_2(\mathbb C)$. Moreover, the first and third printed basis
vectors are reversed relative to the coordinate matrix (1.3). A basis
with respect to which (1.3) \emph{does} represent the adjoint action is
\[
E_1=\begin{pmatrix}0&1\\0&0\end{pmatrix},\qquad
E_2=\frac12\begin{pmatrix}1&0\\0&-1\end{pmatrix},\qquad
E_3=\begin{pmatrix}0&0\\1&0\end{pmatrix}.
\]
Indeed, conjugation by
$\mathrm{diag}(a,d)$ has eigenvalues $a/d,1,d/a$
on this ordered basis, exactly as in (1.3). The basis currently printed
cannot give the stated matrix. \textbf{Action:} correct the descriptive
basis in an erratum and keep the published formula (1.3) only after
checking it against the corrected basis.

\subsection{Two distinct conventions for the outer involution are mixed}
\label{issue:sigma}

\textbf{Location:} \texttt{adjoint\_lift1.tex}, lines 119--131 and
235--255; \texttt{adjoint\_lift4.tex}, lines 50--65;
\textbf{priority: high}. In Section 1, the involution is
\[
\sigma_J(g)=J\,{}^tg^{-1}J,
\qquad J^2=I.
\]
But Lemma 1.10 asserts that $A\sigma([A])$ is represented by
$A\,{}^tA^{-1}$, and Section 4 describes the induced building action as
$[L]\mapsto[A L^\vee]$. Taken literally, the former product for the
stated involution is instead
\[
A J\,{}^tA^{-1}J.
\]
For example, with $A=\mathrm{diag}(a,b,c)$,
$A\,{}^tA^{-1}=I$, whereas
$AJ\,{}^tA^{-1}J=\mathrm{diag}(a/c,1,c/a)$, which
is generally not projectively the identity. Thus the formulas do not
refer to the same representative and convention.

This may be a \emph{repairable normalization issue}, not a defect in the
geometric idea: if $B=AJ$, then
$B\,{}^tB^{-1}=AJ\,{}^tA^{-1}J$. The paper should state the
replacement of the twisted representative $A$ by $AJ$ explicitly at the
point of transition. Note also that $J$ is \emph{involutory}, not
``idempotent'': $J^2=I\ne J$.

\subsection{The centralizer and orbital-integral definitions have wrong
bound variables}
\label{issue:centralizers}

\textbf{Location:} \texttt{adjoint\_lift1.tex}, lines 384--445 and
514--529; \textbf{priority: high}. In (1.15) and (1.16), the element
inside the set-builder is named $g$ or $h$, although $x$ is the element
whose membership in the centralizer is being tested. The intended
formulas are of the type
\[
G(g\sigma,F)=\{x\in G(F):x^{-1}g\sigma(x)=g\},
\qquad
H_i(h,F)=\{x\in H_i(F):x^{-1}hx=h\}.
\]
More seriously, the integrand in (1.17) is printed as
$\phi(g^{-1}g\sigma(g))$: it contains no integration variable.
It must have the form
\[
\phi\bigl(x^{-1}g\sigma(x)\bigr),
\]
as correctly used later in (2.5) and (2.12). In the same passage,
$H_i(h\sigma,F)$ appears where $H_i(h,F)$ is meant.

In (1.20), the twisted double quotient ends with $/H_i(F)$ instead of
$/G(F)$, and the definition of $St_\sigma(g,G)$ uses ordinary
$x^{-1}gx$ rather than $x^{-1}g\sigma(x)$. Its first factor also needs
attention as to whether rational points or geometric points are
intended. These are not innocuous in a definition used to enumerate
rational twisted conjugacy classes. \textbf{Action:} rebuild (1.15)--(1.21)
from a single coherent convention; where appropriate, state the usual
\emph{kernel} of the map in nonabelian $H^1$ rather than an unqualified
$H^1$ identification. In the particular groups at hand a simplification
may be available, but it must be justified.

\subsection{The torus restriction in Section 2 is inconsistent with the
stated $L$-map}
\label{issue:satake}

\textbf{Location:} \texttt{adjoint\_lift2.tex}, lines 185--197,
288--340 and 413--430; \textbf{priority: high}. The text gives
\[
\lambda_2:\mathrm{diag}(t,t^{-1})
\longmapsto\mathrm{diag}(t,t^{-1},1).
\]
Consequently a character represented by the monomial
$z_1^{k_1}z_2^{k_2}z_3^{k_3}$ restricts as
\[
\bigl(z_1^{k_1}z_2^{k_2}z_3^{k_3}\bigr)\circ\lambda_2
=t^{\,k_1-k_2}.
\]
The criterion in the paragraph preceding Lemma 2.11(b), claiming an
``if and only if'' involving
$k_1+k_2=k'_1+k'_2=2k_3=2k'_3$, does not follow
from this map. For instance, the triples $(1,0,0)$ and $(2,1,0)$
both restrict to $t$, although they fail that criterion. Likewise, the
sum in Lemma 2.11(b)(iii), written with $k_1=n_1$ and $k_2=n_2$,
requires a rederivation using the correct quotient of the character
lattice. As presented, the proof of the split $H_2$ identity in (2.13)
is not independently verifiable from the stated lemma.

There is also a repeated root factor in the displayed expression for
$\Delta_G$ immediately before (2.9): it includes
$|1-\gamma_1/\gamma_3|$ twice but omits the root pair $1,2$.
The defining product over all roots is the reliable starting point.
For a diagonal regular element, one can write
\[
\Delta_G(g)=\prod_{1\le i<j\le3}
\left|\left(1-\frac{\gamma_i}{\gamma_j}\right)
\left(1-\frac{\gamma_j}{\gamma_i}\right)\right|^{1/2}.
\]
\textbf{Action:} correct the lattice parametrization and repeat the
coefficient comparison, including the discriminant normalization. The
misstated coefficient criterion is present in the published text too.

\subsection{The averaging subgroup in (3.17)--(3.18) cannot be both compact
and finite index}
\label{issue:U}

\textbf{Location:} \texttt{adjoint\_lift3.tex}, lines 389--412;
\texttt{adjoint\_lift4.tex}, lines 473--480 and 555--561;
\textbf{priority: high}. Section 3 says
$U\subset F^\times$ is a \emph{compact open subgroup of finite index}.
There is no such subgroup in a nonarchimedean local field:
\[
F^\times\simeq\varpi^{\mathbb Z}\times\OF^\times,
\]
and a compact subgroup has valuation image $\{0\}$, hence cannot
have finite index in $F^\times$. Therefore the displayed denominator
$[F^\times:U]$ is not finite under the words as written.

The manuscript itself suggests the repair in (4.15):
\[
U(d)=U_d\varpi^{2\mathbb Z},\qquad
U_d=1+\varpi^d\OF
\quad(\text{for }d\ge1).
\]
Here $U_d$ is compact open in $\OF^\times$, whereas $U(d)$ is open
and of finite index in $F^\times$, but is \emph{not compact}.
For suitable $d$, it lies in the relevant norm subgroup. The statement
of (3.17)--(3.18) should use the finite-index group $U(d)$ or a
precisely specified equivalent, and then prove that the averages are
independent of admissible refinements.

\subsection{Lemma 4.14 is internally contradictory}
\label{issue:parity}

\textbf{Location:} \texttt{adjoint\_lift4.tex}, lines 449--457;
\textbf{priority: high}. The first sentence of Lemma 4.14 states that
the parity of $c$ depends only on the \emph{stable}
$\sigma$-conjugacy class. Its second sentence says that two
stably $\sigma$-conjugate, but not rationally $\sigma$-conjugate,
representatives in the unramified case have \emph{opposite} parities.
These cannot both hold. The preceding paragraph indicates that the first
sentence should likely say \emph{rational $\sigma$-conjugacy class},
while the stable class contains the two possibilities. The same wording
occurs in the 1990 article, p.~1122 \cite{Joyner1990}.

A related problem appears in (4.19),
\texttt{adjoint\_lift4.tex}, lines 645--657: the equality chain begins
with $P'$ and ends in $P$. But $P'$ is an endpoint in $\mathcal Z$,
whereas $P$ is a point with a nondegenerate characteristic segment;
in the argument at hand $P\notin\mathcal Z$. This seems to be an
incorrect equality or an omitted distinction between the endpoint and
its leaf vertex. Because Lemma 4.16 is repeatedly invoked in Section 5,
its proof should be rewritten with the vertices distinctly named.

\subsection{The normalization $Q$ disappears at the decisive point of
Section 5}
\label{issue:Q}

\textbf{Location:} \texttt{adjoint\_lift3.tex}, lines 34--80,
390--405; \texttt{adjoint\_lift5.tex}, lines 26--43 and 312--369;
\textbf{priority: high}. Section 3 sets
\[
Q=(1+q^{-1})(1+q^{-1}+q^{-2}),\qquad Q(0)=Q,
\]
and Lemma 3.3 gives $s(0)=Q$ when the $H_2$ element has one fixed
vertex in the case relevant to the conclusion of Section 5.
The latter explicitly recognizes this fact just before (5.10):
``so $s(0)=Q$.'' Yet the initial case analysis in Section 5 claims
$s(0)=1$, and (5.10) writes $S(0)=1$ after \emph{omitting} the
multiplicative $Q(0)$ of (3.18).

A consistent formulation would be to show that the \emph{normalized
weighted fixed-point sum}, say $S(0)/Q$, equals $1$; the conclusion
would then be
\[
\frac{S(0)}Q=1\quad\Longrightarrow\quad S(0)=Q=s(0).
\]
The equality is arithmetically checkable from the displayed expressions:
the two fixed-point formulas (5.12)--(5.13) differ by $q^m$, as
needed for (5.11), but a factor $Q$ must be restored to obtain the
stated comparison of orbital integrals. This looks like a repairable
normalization omission, not automatically a failure of the count.

\subsection{The ``ramified and close to $-1$'' case really can occur}
\label{issue:ramified}

\textbf{Location:} \texttt{adjoint\_lift5.tex}, lines 34--46
and 371--398; \textbf{priority: high for wording, medium for the proof}.
The announced assertion that an elliptic norm eigenvalue cannot satisfy
both $E/F$ ramified and $|1+\beta_1|_E<1$ is \emph{false literally}.
For an explicit example, take
\[
F=\mathbb Q_3,\qquad E=\mathbb Q_3(\sqrt3),\qquad
\beta=2+\sqrt3.
\]
Then $\beta^{-1}=2-\sqrt3$, so $\Nm_{E/F}(\beta)=1$,
$F(\beta)=E$ is ramified, and
\[
1+\beta=3+\sqrt3,\qquad |1+\beta|_E<1.
\]
The matrix
\[
h=\begin{pmatrix}0&-1\\1&4\end{pmatrix}\in\mathrm{SL}_2(F)
\]
is regular elliptic with these eigenvalues. In particular, this
combination of hypotheses must not simply be declared impossible.

The final paragraph of Section 5 actually studies the case and seeks
to prove that it contributes \emph{zero fixed points} to the relevant
orbital-integral counts. That is the plausible intended statement.
\textbf{Action:} change the case description in a future corrected
edition to ``the ramified case gives no contributing fixed vertices''
(or a precisely stated equivalent), and supply the needed proof. The
example above prevents a misleading literal interpretation.

\subsection{The sign in the last comparison should be proved rather than
inferred from an absolute value}
\label{issue:sign}

\textbf{Location:} \texttt{adjoint\_lift5.tex}, lines 43--49 and
256--310; \textbf{priority: medium/high}. Proposition 5.6 proves only
\[
|S(0)|=s(0)
\quad\text{when }|1+\beta_1|=1,
\]
whereas the desired fundamental identity is $S(0)=s(0)$. The text says
the remaining global sign has been chosen to make $S(0)\ge0$ but also
says this positivity will only be proved later. The signed convention
in (3.18) is tied to a Hilbert symbol, so it cannot be treated as an
independent arbitrary sign at each element. State and verify the sign
normalization in the $|1+\beta_1|=1$ range, perhaps as a short lemma.

\section{Confirmed differences between the uploaded transcription and
the 1990 publication}

A spot comparison is especially useful here because the source file is
dated March 2003 and already advertises minor post-publication edits.
The following differences are visible in the actual published pages,
not merely inferred from mathematical conventions.

\begin{enumerate}[label=\arabic*.]
\item \textbf{The local Langlands sentence was altered incorrectly.}
In \texttt{adjoint\_lift1.tex}, lines 185--190, the sentence begins
``According to the local Langlands correspondence for $GL(3)$'' and
uses a parameter into ${}^L SL(3,\mathbb C)$ while assigning packets
of $SL(2,F)$. The journal version (pp.~1100--1101) says \emph{for
$SL(2)$} and has the parameter into ${}^L SL(2,\mathbb C)$.
The source passage is a transcription error and should be checked
against that page rather than emended by conjecture.

\item \textbf{The target of $\theta^*$ changed in the transcription.}
In \texttt{adjoint\_lift1.tex}, lines 193--197, $\theta^*$ is written
as taking values in ${}^L PGL(2,\mathbb C)$, even though it was
constructed using (1.4). The journal version (p.~1101) has the
target ${}^L PGL(3,\mathbb C)$, consistent with (1.4) and (1.8).

\item \textbf{The Weyl denominator lost a critical ``$1-$''.}
In \texttt{adjoint\_lift3.tex}, equation (3.9), line 223, the term
$1-s_1^{-1}s_2$ has become $s_1^{-1}s_2$. The published version
(p.~1114) includes $1-s_1^{-1}s_2$. The correct displayed product is
\[
\rho(s)\prod_{\alpha>0}(1-\alpha(s)^{-1})
=\frac{s_1}{s_3}
 \left(1-\frac{s_2}{s_1}\right)
 \left(1-\frac{s_3}{s_2}\right)
 \left(1-\frac{s_3}{s_1}\right).
\]
Without that factor, the expression is not alternating in $s_1,s_2,s_3$,
contradicting the sentence immediately following (3.9) and invalidating
the displayed justification of (3.10). This is the most consequential
confirmed error introduced during re-typesetting.
\end{enumerate}

By contrast, the incorrect ``idempotent'' description of $J$, the
conflicting parity sentence of Lemma 4.14, and the $Q$ issue in
Section 5 also appear in the journal version. Their historical
presence should be recorded rather than attributed to the 2003 typing.

\section{Readability and secondary corrections}

\begin{enumerate}[label=\arabic*.]
\item \textbf{State the convention sheet at the outset.} Specify the
normalization of $|\cdot|_F$ and $|\cdot|_E$, the residue cardinality
$q$, the chosen uniformizer, the distinction between $SL_2$, $PGL_2$
and their complex duals, the version of the twist, and the Haar
measure of the hyperspecial compact subgroup. Section 3 switches
between the two meanings of $Q(\Lambda)$ (a double coset and a
scalar factor); reserve different symbols for them.

\item \textbf{Avoid using undeclared $\lambda_1-\lambda_3$.}
The four definitions in \texttt{adjoint\_lift3.tex}, lines 66--73,
contain $q^{\lambda_1-\lambda_3}$, although the associated exponents
were denoted $k_1,k_2,k_3$. Check the intended power against the later
$q^{k_3-k_1}$ and reconcile any sign. This matters for nonunit
characteristic functions, even if the exponent vanishes at $\Lambda=0$.

\item \textbf{Repair ancillary inaccuracies.} In Section 4 the subgroup
$Z(F)$ is described as the center of $PGL_3(F)$; for the quotient
$GL_3(F)/F^\times$ it should refer to the scalar center of $GL_3(F)$.
In Section 1, ``$i=1.2$'' needs a comma, and the formulas distinguishing
ordinary and twisted conjugacy should use different symbols as needed.
The claimed exponential identities in Section 2.2 should be stated as
linearizations (or justified with a triangular Jacobian calculation):
for a nonabelian unipotent group, the exact identity
$nmn^{-1}m^{-1}=\exp((1-\operatorname{Ad}(m))\log n)$ need not hold
without higher commutator terms.

\item \textbf{Make the geometric proof verifiable.} Lemmas 4.9, 4.13,
4.16, 4.18, 5.3 and 5.4 are invoked at crucial points, but several
proofs are just sketches or end with ``similar'' and ``easy''.
In a modern explanatory edition, add explicit descriptions of the
lattice coordinates and the action of $A_a\sigma$; prove the orbit and
stabilizer counts needed for (4.16)(e). In Lemma 5.4, the valuation
analysis for $b_1-b_2=-1,0$ should not be omitted.

\item \textbf{Clarify the statement actually proved.} The introduction
mentions arbitrary matching spherical functions, but the principal
original calculation in Section 5 treats the \emph{unit} element of the
Hecke algebra and assumes odd residual characteristic. Distinguish
the previously known full fundamental lemma (attributed to Langlands)
from the case directly demonstrated here. The introductory discussion
of the conjectural transfer should also be set explicitly in its
historical context rather than presented as the present state of the
subject.

\item \textbf{Presentation.} The source compiles successfully as
12-point Times-style \LaTeX{} on US letter paper with resolved cross
references. There remain underfull/overfull boxes and a warning about
the older \verb|\over| form. The four figure environments in
\texttt{adjoint\_lift4.tex} contain the drawings necessary to follow
the paper; the stand-alone \texttt{adjoint\_lift\_pics.tex} repeats
some of the same graphics and should not be compiled as an additional
copy of those figures in the main text. Consider enlarging figure labels,
renaming ``characteristic leafs'' as ``characteristic leaves'', and
replacing raw picture coordinates with a preserved, separately
verifiable vector source in a future edition.
\end{enumerate}

\section{Verification checklist and proposed editorial route}

Before an unqualified corrected edition is released, I recommend the
following sequence. It deliberately separates transcription fidelity
from mathematical correction, respecting the historical nature of the
paper.

\begin{enumerate}[label=\textbf{Step \arabic*.},leftmargin=*]
\item \textbf{Create a fidelity log}: systematically compare
(1.3), (1.15)--(1.21), (2.7)--(2.13), (3.9)--(3.18), Lemma 4.14,
(4.19), and (5.10)--(5.13) against the published pages.
Mark which deviations are 2003 transcription errors and which were
present in 1990.
\item \textbf{Create a separate mathematical errata appendix}: address
the twist convention, centralizer definitions, Satake restriction,
averaging subgroup and normalization. Preserve the historical text if
an archival transcription is the goal, and link each erratum to the
relevant numbered assertion instead of altering it silently.
\item \textbf{Re-derive the nonunit-to-unit reduction}: verify the
$\lambda_i^*$ coefficient identities and the $Q(\Lambda)$ measure
factors independently of the manuscript's formulas. Confirm that the
fixed-point averages lead to exactly $R(0)=r(0)$ and $S(0)=s(0)$,
with no lost factor $Q$ or undetermined sign.
\item \textbf{Strengthen the Section 5 case split}: explicitly allow
ramified $\beta$ close to $-1$ and prove the actual vanishing statement;
record all hypotheses and justify the signed comparison.
\item \textbf{Polish the explanatory layer last}: add a compact glossary,
a paragraph introducing buildings and dual lattice classes, precise
captions, and a short account of what is original in Section 5.
\end{enumerate}

\section{Concluding assessment}

The paper's general method and historical place remain of clear
interest. I did not find a counterexample to the intended equality of
the \emph{properly normalized} orbital integrals. I did find
counterexamples or direct contradictions to \emph{some of the statements
as printed}, including the finite-index/compactness claim for $U$,
the two sentences of Lemma 4.14, and the literal assertion that the
ramified near-$-1$ case cannot occur. There is also a missing factor
$Q$ at the final normalization step. Taken together, these warrant a
careful correction and independent check of the delicate steps before
the uploaded version is used as a reliable standalone proof.

For an archival edition, retain the original article text intact and
publish this report or a more formal \emph{Editor's Notes and Errata}
appendix alongside it. The confirmed 2003 transcription mistakes should
be documented separately from questions that already occur in the 1990
publication.

\begin{thebibliography}{9}
\bibitem{Joyner1990}
D.~Joyner, ``On twisted orbital integral identities for $PGL(3)$ over a
$p$-adic field,'' \emph{Canadian Journal of Mathematics}
\textbf{42} (1990), no.~6, 1098--1130.
\url{https://doi.org/10.4153/CJM-1990-059-6}.
\end{thebibliography}

\end{document}
